\begin{abstract}
Attention--FC 分离式（Attention-FC Disaggregated，AFD）大语言模型推理系统将受内存带宽约束的 Attention 操作卸载到内存资源丰富的加速器（如 CPU、HBM-PIM），同时把受计算约束的全连接（Fully-Connected，FC）操作保留在 GPU 上。
本文首先设计分离式屋顶线模型（Disaggregated Roofline Model，DRM）来刻画 AFD 的性能；该模型揭示，\emph{系统吞吐量受加速器的限制因素约束，而这一因素可能是内存带宽，也可能是内存容量}。
我们观察到，以往 AFD 系统往往忽视且未能平衡这些约束，因而导致资源利用不足或吞吐量受限。
为此，我们提出 \sys：首个集成 \textit{DIMM-PIM} 的 AFD 系统；DIMM-PIM 是一类凭借可扩展容量与带宽实现二者平衡的新型加速器。
为解决 DIMM-PIM 中分布式协作 DRAM 芯片固有的同步难题，\sys 采用\emph{无气泡流水线}和\emph{混合粒度重布局}来高效计算 Attention。
此外，它通过\emph{秩组粒度的通信--计算重叠}与\emph{对齐预测调度}，最大化跨设备资源利用率。
评估结果表明，与最先进的 HBM-PIM 方案相比，\sys 可实现最高 5.15$\times$ 的加速。
\end{abstract}
